% MATH3701 Supplementary Proofs — shared preamble
% % MATH3701 Supplementary Proofs — shared preamble
% % MATH3701 Supplementary Proofs — shared preamble
% \input{proof_preamble} at the top of each proof file

\documentclass[11pt,a4paper]{article}
\usepackage{amsmath,amssymb,amsthm}
\usepackage{mathtools}
\usepackage{bm}
\usepackage{geometry}
\geometry{margin=2.5cm}
\usepackage{hyperref}
\hypersetup{colorlinks=true, linkcolor=blue, urlcolor=blue}

% ---- Math shortcuts (matching slides/preamble.tex) ----
\newcommand{\E}{\mathbb{E}}
\newcommand{\Var}{\mathrm{Var}}
\newcommand{\Cov}{\mathrm{Cov}}
\newcommand{\Prob}{\mathbb{P}}
\newcommand{\R}{\mathbb{R}}
\newcommand{\N}{\mathcal{N}}
\newcommand{\bX}{\mathbf{X}}
\newcommand{\bx}{\mathbf{x}}
\newcommand{\by}{\mathbf{y}}
\newcommand{\bY}{\mathbf{Y}}
\newcommand{\bb}{\mathbf{b}}
\newcommand{\ba}{\mathbf{a}}
\newcommand{\bbeta}{\boldsymbol{\beta}}
\newcommand{\beps}{\boldsymbol{\epsilon}}
\newcommand{\bmu}{\boldsymbol{\mu}}
\newcommand{\btheta}{\boldsymbol{\theta}}
\newcommand{\bU}{\mathbf{U}}
\newcommand{\bJ}{\mathbf{J}}
\newcommand{\bI}{\mathbf{I}}
\newcommand{\bH}{\mathbf{H}}
\newcommand{\bK}{\mathbf{K}}
\newcommand{\bL}{\mathbf{L}}
\newcommand{\bM}{\mathbf{M}}
\newcommand{\bS}{\mathbf{S}}
\newcommand{\bW}{\mathbf{W}}
\newcommand{\bz}{\mathbf{z}}
\newcommand{\boldeta}{\boldsymbol{\eta}}
\newcommand{\logit}{\mathrm{logit}}
\newcommand{\iid}{\stackrel{\mathrm{iid}}{\sim}}
\newcommand{\tr}{\mathrm{tr}}

% ---- Theorem environments ----
\theoremstyle{plain}
\newtheorem{proposition}{Proposition}
\newtheorem{lemma}{Lemma}
\newtheorem{theorem}{Theorem}
\theoremstyle{definition}
\newtheorem{remark}{Remark}

% ---- Module branding ----
\newcommand{\moduleheader}{%
  \noindent\textbf{MATH3701 Statistical Modelling}\\
  School of Mathematics, University of Leeds\\[6pt]
}
 at the top of each proof file

\documentclass[11pt,a4paper]{article}
\usepackage{amsmath,amssymb,amsthm}
\usepackage{mathtools}
\usepackage{bm}
\usepackage{geometry}
\geometry{margin=2.5cm}
\usepackage{hyperref}
\hypersetup{colorlinks=true, linkcolor=blue, urlcolor=blue}

% ---- Math shortcuts (matching slides/preamble.tex) ----
\newcommand{\E}{\mathbb{E}}
\newcommand{\Var}{\mathrm{Var}}
\newcommand{\Cov}{\mathrm{Cov}}
\newcommand{\Prob}{\mathbb{P}}
\newcommand{\R}{\mathbb{R}}
\newcommand{\N}{\mathcal{N}}
\newcommand{\bX}{\mathbf{X}}
\newcommand{\bx}{\mathbf{x}}
\newcommand{\by}{\mathbf{y}}
\newcommand{\bY}{\mathbf{Y}}
\newcommand{\bb}{\mathbf{b}}
\newcommand{\ba}{\mathbf{a}}
\newcommand{\bbeta}{\boldsymbol{\beta}}
\newcommand{\beps}{\boldsymbol{\epsilon}}
\newcommand{\bmu}{\boldsymbol{\mu}}
\newcommand{\btheta}{\boldsymbol{\theta}}
\newcommand{\bU}{\mathbf{U}}
\newcommand{\bJ}{\mathbf{J}}
\newcommand{\bI}{\mathbf{I}}
\newcommand{\bH}{\mathbf{H}}
\newcommand{\bK}{\mathbf{K}}
\newcommand{\bL}{\mathbf{L}}
\newcommand{\bM}{\mathbf{M}}
\newcommand{\bS}{\mathbf{S}}
\newcommand{\bW}{\mathbf{W}}
\newcommand{\bz}{\mathbf{z}}
\newcommand{\boldeta}{\boldsymbol{\eta}}
\newcommand{\logit}{\mathrm{logit}}
\newcommand{\iid}{\stackrel{\mathrm{iid}}{\sim}}
\newcommand{\tr}{\mathrm{tr}}

% ---- Theorem environments ----
\theoremstyle{plain}
\newtheorem{proposition}{Proposition}
\newtheorem{lemma}{Lemma}
\newtheorem{theorem}{Theorem}
\theoremstyle{definition}
\newtheorem{remark}{Remark}

% ---- Module branding ----
\newcommand{\moduleheader}{%
  \noindent\textbf{MATH3701 Statistical Modelling}\\
  School of Mathematics, University of Leeds\\[6pt]
}
 at the top of each proof file

\documentclass[11pt,a4paper]{article}
\usepackage{amsmath,amssymb,amsthm}
\usepackage{mathtools}
\usepackage{bm}
\usepackage{geometry}
\geometry{margin=2.5cm}
\usepackage{hyperref}
\hypersetup{colorlinks=true, linkcolor=blue, urlcolor=blue}

% ---- Math shortcuts (matching slides/preamble.tex) ----
\newcommand{\E}{\mathbb{E}}
\newcommand{\Var}{\mathrm{Var}}
\newcommand{\Cov}{\mathrm{Cov}}
\newcommand{\Prob}{\mathbb{P}}
\newcommand{\R}{\mathbb{R}}
\newcommand{\N}{\mathcal{N}}
\newcommand{\bX}{\mathbf{X}}
\newcommand{\bx}{\mathbf{x}}
\newcommand{\by}{\mathbf{y}}
\newcommand{\bY}{\mathbf{Y}}
\newcommand{\bb}{\mathbf{b}}
\newcommand{\ba}{\mathbf{a}}
\newcommand{\bbeta}{\boldsymbol{\beta}}
\newcommand{\beps}{\boldsymbol{\epsilon}}
\newcommand{\bmu}{\boldsymbol{\mu}}
\newcommand{\btheta}{\boldsymbol{\theta}}
\newcommand{\bU}{\mathbf{U}}
\newcommand{\bJ}{\mathbf{J}}
\newcommand{\bI}{\mathbf{I}}
\newcommand{\bH}{\mathbf{H}}
\newcommand{\bK}{\mathbf{K}}
\newcommand{\bL}{\mathbf{L}}
\newcommand{\bM}{\mathbf{M}}
\newcommand{\bS}{\mathbf{S}}
\newcommand{\bW}{\mathbf{W}}
\newcommand{\bz}{\mathbf{z}}
\newcommand{\boldeta}{\boldsymbol{\eta}}
\newcommand{\logit}{\mathrm{logit}}
\newcommand{\iid}{\stackrel{\mathrm{iid}}{\sim}}
\newcommand{\tr}{\mathrm{tr}}

% ---- Theorem environments ----
\theoremstyle{plain}
\newtheorem{proposition}{Proposition}
\newtheorem{lemma}{Lemma}
\newtheorem{theorem}{Theorem}
\theoremstyle{definition}
\newtheorem{remark}{Remark}

% ---- Module branding ----
\newcommand{\moduleheader}{%
  \noindent\textbf{MATH3701 Statistical Modelling}\\
  School of Mathematics, University of Leeds\\[6pt]
}


\begin{document}

\moduleheader

\begin{center}
  {\Large\bfseries Supplementary Proof: Optimality of Natural Splines}
\end{center}

\bigskip

\noindent This note proves that among all sufficiently smooth functions
interpolating prescribed values at a set of knots, the natural cubic
spline has the smallest roughness.  This result underpins Chapter~11
of the notes (Propositions 11.1 and 11.2).

\section*{Setup and statement}

Let $t_1 < t_2 < \cdots < t_n$ be $n$ distinct knots and let
$y_1,\dots,y_n$ be prescribed values.  We consider functions $f$ in
the Sobolev space $W^{2,2}(\R)$: functions that are twice
differentiable with square-integrable second derivative.  The
\emph{roughness} of $f$ is
\[
  J(f) = \int_{-\infty}^{\infty} [f''(t)]^2\,dt.
\]

\begin{theorem}[Optimality of the natural cubic spline]
\label{thm:optimality}
  Let $f^*$ be the natural cubic spline with knots at $t_1,\dots,t_n$
  satisfying $f^*(t_i) = y_i$ for $i=1,\dots,n$.  Then, for any other
  function $g \in W^{2,2}(\R)$ satisfying $g(t_i)=y_i$,
  \[
    J(g) \geq J(f^*),
  \]
  with equality if and only if $g = f^*$.
\end{theorem}

\section*{Key property of the natural cubic spline}

The proof rests on the following orthogonality lemma.

\begin{lemma}
\label{lem:orthog}
  Let $f^*$ be the natural cubic spline with $f^*(t_i)=y_i$.  For any
  $h\in W^{2,2}(\R)$ with $h(t_i) = 0$ for all $i$,
  \[
    \int_{-\infty}^{\infty} (f^*)''(t)\,h''(t)\,dt = 0.
  \]
\end{lemma}

\begin{proof}
  Since $f^*$ is a natural cubic spline, it is a piecewise cubic
  polynomial with the natural boundary conditions
  $(f^*)''(t) \to 0$ as $t\to\pm\infty$
  (in fact $(f^*)'' = 0$ outside $[t_1, t_n]$).
  Integrate by parts twice:
  \begin{align*}
    \int_{-\infty}^{\infty} (f^*)''(t)\,h''(t)\,dt
    &= \bigl[(f^*)''(t)\,h'(t)\bigr]_{-\infty}^{\infty}
       - \int_{-\infty}^{\infty}(f^*)'''(t)\,h'(t)\,dt.
  \end{align*}
  The boundary term vanishes because $(f^*)'' = 0$ outside the knot
  range, so $(f^*)''(\pm\infty)=0$.  Integrate by parts again:
  \begin{align*}
    -\int_{-\infty}^{\infty}(f^*)'''(t)\,h'(t)\,dt
    &= -\bigl[(f^*)'''(t)\,h(t)\bigr]_{-\infty}^{\infty}
       + \int_{-\infty}^{\infty}(f^*)''''(t)\,h(t)\,dt.
  \end{align*}
  Since $f^*$ is linear outside the knot range, $(f^*)'''=0$ there,
  so the boundary term vanishes.

  Inside the knot range, a cubic spline has $(f^*)''''(t)=0$
  on each open interval $(t_{i-1},t_i)$.  Thus the integral reduces to
  a sum of contributions at the knots.  Specifically, if we write the
  fourth derivative in the distributional sense, the jumps in
  $(f^*)'''$ at each knot $t_i$ contribute a Dirac mass:
  \[
    \int_{-\infty}^{\infty}(f^*)''''(t)\,h(t)\,dt
    = \sum_{i=1}^n \bigl[(f^*)'''(t_i^+) - (f^*)'''(t_i^-)\bigr]h(t_i)
    = 0,
  \]
  since $h(t_i) = 0$ for all $i$.  Hence
  $\int (f^*)'' h'' \,dt = 0$.
\end{proof}

\section*{Proof of Theorem \ref{thm:optimality}}

Let $g$ be any function in $W^{2,2}(\R)$ with $g(t_i)=y_i$, and set
$h = g - f^*$.  Then $h(t_i) = 0$ for all $i$.  We decompose:
\begin{align*}
  J(g)
  &= \int_{-\infty}^{\infty} [g''(t)]^2\,dt
   = \int_{-\infty}^{\infty} \bigl[(f^*)''(t) + h''(t)\bigr]^2\,dt \\
  &= J(f^*) + 2\int_{-\infty}^{\infty}(f^*)''(t)\,h''(t)\,dt + J(h).
\end{align*}
By Lemma~\ref{lem:orthog}, the cross-term vanishes, giving
\[
  J(g) = J(f^*) + J(h) \geq J(f^*),
\]
since $J(h) = \int[h'']^2\,dt \geq 0$.

Equality holds if and only if $J(h)=0$, i.e., $h''=0$ everywhere,
meaning $h$ is a linear function.  If $h(t_i)=0$ for all $i\geq 2$
distinct knots, then $h\equiv 0$.  Hence $g = f^*$, proving
uniqueness. \hfill$\square$

\section*{Representations of natural splines (Proposition~11.1)}

The proof above also establishes that the natural cubic spline in
$W^{2,2}(\R)$ interpolating $n$ points can be represented in the form
\[
  f^*(t) = a_0 + a_1 t + \sum_{i=1}^n b_i\,|t-t_i|^3,
\]
subject to the constraints $\sum b_i = 0$ and $\sum b_i t_i = 0$
(which ensure $f''=0$ outside $[t_1,t_n]$).  This gives
$n+2$ parameters with $2$ constraints, leaving $n$ free parameters
to interpolate the $n$ data values.

For the linear case ($\nu=1$), the analogous representation is
\[
  f^*(t) = a_0 + \sum_{i=1}^n b_i\,|t-t_i|,
\]
subject to $\sum b_i = 0$, and the optimality proof follows the same
argument using $J(f) = \int [f'(t)]^2 dt$ and a single integration
by parts.

\end{document}
